\documentclass{article}

\begin{document}

\section{Context}

This fictional study considers how a researcher can find notes from a completed
interview study. The notes remain in ordinary text files, but the reason for a
particular design decision can be hard to find when related passages are spread
across several files.

The project explores a small local index that helps the researcher find those
passages without moving the notes into a separate online service. The index is
an aid to navigation; the original notes remain the source of the evidence.

\section{Prototype}

The prototype reads a selected folder and stores a searchable index of the text.
Each result points back to the file and passage where the matching words appear.
The researcher can then read the surrounding notes before interpreting a match.

The search was making it easier to find the notes, and it was useful for
helping the researcher to find things. The prototype presents a match as
\emph{a starting point for inspection}, rather than a conclusion about what
the interview participant meant.

\section{Evaluation}\label{sec:evaluation}
For this example, the researcher prepared \(n = 12\) searches against a
fictional set of notes. Before running a search, they identified the passage
that should answer it. The evaluation records whether that passage appears in
the results, not whether the interpretation of the passage is correct.

The match rate is \(r = m/n\), where \(m\) is the number of searches that
return their prepared target. As described in Section~\ref{sec:evaluation},
this measure applies to the prepared searches and does not establish that the
prototype improves every research workflow.

\end{document}
